\subsection{The remaining two-object budgets after combined source reductions}
\label{CP2:sec-combined-small-corners}

\begin{proposition}[Proved: a reduced finite scope for the six remaining budgets]
\label{CP2:A-combined-small-corner-scope}
Use the jointly oriented source convention of
Corollary~\ref{CP2:A-two-object-eighteen-oriented-twenty-two} and the already proved
exclusion through seventeen arrows.  The band-source portions of the six
remaining eighteen-arrow budgets have the following complete source partition.
The reductions are applied in the displayed order.  The last column counts the source classes retained by precisely these
three source cuts. Later independent source reductions are stated separately;
this table asserts no completion of a finite action computation.
\[
\begin{array}{c|r|r|r|r|r}
(|M|,|N|,|S|)&\text{all}&\text{graph}&\text{semicentral}&
 \text{integral pair}&\text{after three cuts}\\\hline
(6,4,8)&5&0&1&1&3\\
(7,3,8)&43&28&1&5&9\\
(8,3,7)&510&252&38&82&138\\
(9,3,6)&8044&3392&897&1560&2195\\
(10,3,5)&157131&62093&21946&31760&41332\\
(10,4,4)&157131&62093&21946&31760&41332
\end{array}
\]
For $(6,4,8)$, only the nilpotent four-element opposite corner requires
finite action computations.  For $(10,4,4)$, only the five other
nonprojective four-element opposite corners require such computations.
The separately certified nontrivial-minimal-group source pieces of orders
nine and ten remain required when joining a whole budget.
\end{proposition}

\begin{proof}
The graph-source column uses
Theorem~\ref{CP2:A-projective-local-row-graph-obstruction}.  For an opposite
three-element corner there is exactly one nonprojective type, and its
augmented action complex has the required graph form.  For a four-element
corner, the graph comparison covers five of the six nonprojective types.
When $|S|\le4$, the remaining nilpotent type is excluded by
Proposition~\ref{CP2:A-nilpotent-four-corner-four-states}.  Thus the same
graph-eligible source classes may be omitted altogether for $(10,4,4)$.
For $(6,4,8)$ all five source classes are graph-eligible, but the nilpotent
opposite corner must still be retained.

The next column uses a proper idempotent $e$ with $Me=eMe$.  Adjoining the
retract object and deleting the old source replaces its corner by $eMe$ and
its connecting act by $Se$, preserving the actual homotopy type and reducing
the arrow count.  The integral-pair column then uses distinct idempotents
$a,e$ with $ae=ea=e$, the congruence identifying only $a$ and $e$, and the
central integral idempotent $a-e$.  The resulting actual category quotient
is an integral homology equivalence and is surjective on fundamental groups.
For a simply connected $\CP^2$ model it is therefore a homotopy equivalence
to a category with fewer arrows.  These are precisely the previously proved
semicentral and integral-pair reduction theorems
(\ref{CP2:A-extreme-idempotent-corner-reduction} and
\ref{CP2:A-idempotent-pair-integral-reduction}).  They are used only after
the earlier columns, so overlaps are not counted twice.

All sets in the table are joined to the complete literal positive source
catalogs.  The independent source scans test all $165734$ tables in orders
five through ten.  The semicentral test compares the two literal sets
$Me$ and $eMe$.  The integral-pair test checks the sole-pair congruence and
each integral coefficient in $(a-e)x=x(a-e)$.  The second implementation
uses signed coefficient vectors rather than the producer's unordered
two-term multiset equality.  The joins verify the exact source-table and
graph-flag hashes, then take set differences in the order displayed above.
For order seven the nine remaining indices are
$3,15,19,20,28,34,39,41,42$.

At order six the surviving source indices are $0,2,4$.  For the two types
whose outside element is an involution, a row-independent right band act is
excluded by the proved permutation-source sign lemma
\ref{CP2:A-row-independent-permutation-source}: its row character is
projective over $\mathbb Q[C_2]$, and the opposite corner has order four.
In a remaining action of the nilpotent corner, let $F$ denote the image of
its unique nonidempotent element.  If $F^2=F$, its defining relations imply
$F=P_0$.  The arbitrary-state idempotent-action theorem
\ref{CP2:A-nilpotent-four-idempotent-action} then gives vanishing
degree-two products.  Hence the finite computation may restrict to
$F^2\ne F$.  This condition is invariant under conjugation of the state set.
No unproved graph comparison for the nilpotent corner is used.
\end{proof}

\paragraph{Later source reductions.}
Theorems~\ref{CP2:A-seven-source-three-right-zero} and
\ref{CP2:A-eight-source-five-right-zero} separately remove source $3$
at order seven and source $5$ at order eight, respectively, opposite $R_3$.
Their statements have no state-count bound. They are not silently folded into
the three-cut table or its frozen control streams. A successor numerical
scope must bind the literal recognition predicate and its own control before
using either additional omission.

\paragraph{Executed controls and the completion boundary.}
The separately named combined producer and independent scope checker are
exact seven-edit and nine-edit derivatives of the previously checked engines.
The producer propagates all multiplication constraints; the scope checker
assigns actors in fixed order and independently reconstructs the state
centralizers, stabilizers, and orbit weights.  Both have received full source
reviews.  Four actual compilation processes exited zero, and the two builds
of each executable were byte-identical.

Three complete controls compare with earlier closed literal streams:
\[
\begin{array}{c|r|r|r|r}
(|M|,|N|,|S|)&\substack{\text{old right}\\\text{orbits}}&\substack{\text{retained right}\\\text{orbits}}&
 \substack{\text{retained biset}\\\text{orbits}}&\substack{\text{literal zero}\\\text{squares}}\\\hline
(7,3,7)&5963&3791&99336&2458\\
(6,4,7)&326&118&6485&354\\
(9,4,4)&199316&55837&460303&48167
\end{array}
\]
Every surviving action record and literal certificate agrees byte for byte
after right-orbit renumbering.  Every surviving scope count is reconstructed
from the old records and their stabilizer weights.  The fixed-order search
node counts agree as well, except in the $106$ nilpotent-actor scopes where
the explicitly restricted search tree is checked to be no larger.  All
producer, scope, literal, and stream-link processes exited zero.  For the
middle control a fresh original-engine baseline was first completed and
independently checked, because the historical seventeen-arrow stream preceded
the full proper-biset-retraction gate.  That baseline's four stages also
exited zero.

The closed control receipts are
\path{combined_small_corners/control_m7_n3_s7.complete_verified.json},
\path{control_m6_n4_s7.complete_verified.json}, and
\path{control_m9_n4_s4.complete_verified.json} under the registered
two-object working directory.  These controls validate the new finite
scope; they do not themselves close any of the six eighteen-arrow budgets.
Each production profile needs its own closed independent scope, literal
certificate, and exact stream-link receipts before the global inventory can
advance beyond sixteen of twenty-two budgets.
